\documentclass[11pt,a4paper]{article}
\usepackage[margin=24mm]{geometry}
\usepackage{amsmath,amssymb,amsthm}
\title{Class Field Theory I.3.6}
\author{J. S. Milne}
\date{Fixed author PDF v4.03, 2020-08-06}
\begin{document}
\maketitle

\section*{Original statement}
\noindent\emph{Source \texttt{CFT.pdf}, physical page 47.}\par\medskip
% BEGIN CFT_I36_STATEMENT
\noindent\textsc{Theorem 3.6}
{\itshape Let $K_{\pi,n}$ denote the subfield $K[\Lambda_n]$ of
$K^{\mathrm{al}}$ generated over $K$ by the elements of $\Lambda_n$.
\begin{enumerate}
\renewcommand{\labelenumi}{(\alph{enumi})}
\item For each $n$, $K_{\pi,n}/K$ is totally ramified of degree $(q-1)q^{n-1}$.
\item The action of $A$ on $\Lambda_n$ defines an isomorphism
\[
(A/\mathfrak m^n)^\times\longrightarrow\operatorname{Gal}(K_{\pi,n}/K).
\]
In particular, $K_{\pi,n}/K$ is abelian.
\item For each $n$, $\pi$ is a norm from $K_{\pi,n}$.
\end{enumerate}}
% END CFT_I36_STATEMENT

\section*{Original proof}
\noindent\emph{Source \texttt{CFT.pdf}, physical pages 47--48.}\par\medskip
% BEGIN CFT_I36_PROOF_47
\noindent\textsc{Proof.} Again, we may assume that $f\in\mathcal F_\pi$
is a polynomial of the form $\pi T+\cdots+T^q$.

(a), (b). Choose a nonzero root $\pi_1$ of $f(T)$ and (inductively) a root
$\pi_n$ of $f(T)-\pi_{n-1}$. Consider the sequence of fields
\[
K[\Lambda_n]\supset K[\pi_n]\overset{q}{\supset}K[\pi_{n-1}]
 \overset{q}{\supset}\cdots\overset{q}{\supset}K[\pi_1]
 \overset{q-1}{\supset}K.
\]
% END CFT_I36_PROOF_47
% BEGIN CFT_I36_PROOF_48
\noindent Each extension is Eisenstein (ANT, 7.55) with the degree indicated.
Therefore $K[\pi_n]$ is totally ramified over $K$ of degree $q^{n-1}(q-1)$.

Recall that $\Lambda_n$ is the set of roots of $f^{(n)}$ in $K^{\mathrm{al}}$,
and so $K[\Lambda_n]$ is the splitting field of $f^{(n)}$. Therefore
$\operatorname{Gal}(K[\Lambda_n]/K)$ can be identified with a subgroup of
the group of permutations of the set $\Lambda_n$, but Lemma 3.5 implies
that each element of $\operatorname{Gal}(K[\Lambda_n]/K)$ acts on $\Lambda_n$
as an $A$-module isomorphism, and so the image of
$\operatorname{Gal}(K[\Lambda_n]/K)$ in the automorphism group of $\Lambda_n$
is contained in
\[
\operatorname{End}_A(\Lambda_n)=(A/(\pi^n))^\times.
\]
Hence
\[
(q-1)q^{n-1}\geq
 |\operatorname{Gal}(K[\Lambda_n]/K)|=[K[\Lambda_n]:K]
 \geq[K[\pi_n]:K]=(q-1)q^{n-1}.
\]
We must have equalities throughout, and so
$\operatorname{Gal}(K[\Lambda_n]/K)\simeq(A/\mathfrak m^n)^\times$
and $K[\Lambda_n]=K[\pi_n]$.

\clearpage
\noindent(c) Let $f^{[n]}(T)=(f/T)\circ f\circ\cdots\circ f$ ($n$ terms), so that
\[
f^{[n]}(T)=\pi+\cdots+T^{(q-1)q^{n-1}}.
\]
Then $f^{[n]}(\pi_n)=f^{[n-1]}(\pi_{n-1})=\cdots=f(\pi_1)=0$.
Because $f^{[n]}$ is monic of degree $(q-1)q^{n-1}=[K[\pi_n]:K]$,
it must be the minimal polynomial of $\pi_n$ over $K$. Therefore,
\[
\begin{aligned}
\operatorname{Nm}_{K[\Lambda_n]/K}\pi_n&=(-1)^{(q-1)q^{n-1}}\pi\\
 &=\pi\quad\text{unless $q=2$ and $n=1$.}
\end{aligned}
\]
In the exceptional case, $K[\Lambda_1]=K$, and $\pi$ is certainly a norm.
\qed
% END CFT_I36_PROOF_48

\section*{Published mathematical erratum}
\noindent\emph{Milne's archived CFT errata, under \textquotedblleft From Raghuram
Sundararajan\textquotedblright. The original proof above is reproduced as printed.}
\par\medskip
% BEGIN CFT_I36_ERRATUM
\noindent (mathematical) In the proof of Theorem 3.6 (c) (page 39), that the
uniformiser $\pi$ is a norm from $K_{\pi,n}$: \textquotedblleft unless $q=2$ and $n=1$\textquotedblright{}
should be \textquotedblleft unless $q$ is even/(a power of $2$) and $n=1$\textquotedblright.
The argument given to cover that case then also needs to be rectified; this can be done
by taking the norm of $-\pi_1$ rather than the norm of $\pi_1$. (More
precisely, for $q$ a power of $2$,
$\operatorname{Nm}_{K_{\pi,n}/K_{\pi,n-1}}(\pm\pi_n)=\pi_{n-1}$
for $n\geq 2$, while
$\operatorname{Nm}_{K_{\pi,1}/K}(\pm\pi_1)=\mp\pi$, so that $-\pi_n$
works for all $n$. In fact, $\operatorname{Nm}_{K_{\pi,n}/K}(-\pi_n)=\pi$
for any $q$ and $n$, so that the proof can be given without any separate
edge cases.)
% END CFT_I36_ERRATUM

\clearpage
\section*{Original supporting excerpts}
\noindent\emph{Source \texttt{CFT.pdf}, physical page 46, continuing from page 45.}\par\smallskip
\begingroup
\setlength{\abovedisplayskip}{8pt plus 2pt minus 2pt}
\setlength{\belowdisplayskip}{8pt plus 2pt minus 2pt}
% BEGIN CFT_I36_CONTEXT_46
\noindent to take $f=\pi T+T^q$.\footnote[2]{Alternatively, the $p$-adic
Weierstrass preparation theorem (Washington 1997, 7.3) implies that every
$f\in\mathcal F_\pi$ factors into $f_1(T)u(T)$, where $f_1(T)$ is a
polynomial $\pi T+\cdots+aT^q$, $a\equiv 1\bmod\mathfrak m$, and $u(T)$
is a unit in $A[[T]]$. Obviously, the roots of $f^{(n)}$ are the roots of
$f_1^{(n)}$.} Then,
\[
\begin{aligned}
(f\circ f)(T)&\overset{\mathrm{def}}{=}f(f(T))
 =\pi(\pi T+\cdots+T^q)+\cdots+(\pi T+\cdots+T^q)^q\\
 &=\pi^2T+\cdots+T^{q^2},
\end{aligned}
\]
and
\[
f^{(n)}(T)=\pi^nT+\cdots+T^{q^n}.
\]
From the Newton polygon of $f^{(n)}$ (cf. ANT, 7.44), we see that the
roots of $f^{(n)}$ all have positive $\operatorname{ord}_K$, hence absolute
value $<1$, and so $\Lambda_n$ is the set of \emph{all} roots of $f^{(n)}$
in $K^{\mathrm{al}}$ endowed with the commutative group stucture
\[
\alpha+_{F_f}\beta=F_f(\alpha,\beta)=\alpha+\beta+\cdots,
\]
and the $A$-module structure,
\[
[a]_f\alpha=a\alpha+\cdots.
\]

\noindent\textsc{Example 3.2} Take $K=\mathbb Q_p$ and
$f=(T+1)^p-1\in\mathcal F_p$; then
\[
\Lambda_n=\{\alpha\in\mathbb Q_p^{\mathrm{al}}\mid(\alpha-1)^{p^n}=1\}
 \simeq\{\zeta\mid\zeta^{p^n}=1\}=\mu_{p^n}.
\]
The addition $+_F$ on $\Lambda_n$ corresponds to multiplication on
$\mu_{p^n}$, and the $\mathbb Z_p$-module structure is as defined before.
It follows that $\Lambda_n\approx\mathbb Z/p^n\mathbb Z$ (as a
$\mathbb Z_p$-module).

Recall that a \emph{\textbf{cyclic module}} is a module that can be
generated by a single element. Because $A$ is a principal ideal domain
with only one prime element up to conjugates, every finitely generated
torsion $A$-module $M$ decomposes into a direct sum of cyclic modules
\[
M\approx A/(\pi^{n_1})\oplus\cdots\oplus A/(\pi^{n_r}),
\qquad n_1\leq n_2\leq\cdots,
\]
and the sequence $n_1,\ldots,n_r$ is uniquely determined.

\noindent\textsc{Lemma 3.3}
{\itshape Let $M$ be an $A$-module, and let
$M_n=\operatorname{Ker}(\pi^n\colon M\longrightarrow M)$. Assume,
\begin{enumerate}
\renewcommand{\labelenumi}{(\alph{enumi})}
\item $M_1$ has $q\overset{\mathrm{def}}{=}(A:(\pi))$ elements, and
\item $\pi\colon M\longrightarrow M$ is surjective.
\end{enumerate}
Then $M_n\approx A/(\pi^n)$; in particular, it has $q^n$ elements.}
\par\medskip
\noindent\textsc{Proof.} We use induction on $n$. Because $A/\pi$ has
order $q$, condition (a) and the structure theorem imply that
$M_1\approx A/(\pi)$. Consider the sequence
\[
0\longrightarrow M_1\longrightarrow M_n\xrightarrow{\pi}M_{n-1}
 \longrightarrow 0.
\]
Condition (b) implies that it is exact at $M_{n-1}$, and is therefore
exact. It follows that $M_n$ has $q^n$ elements. If $M_n$ is not cyclic,
\[
M_n\approx A/(\pi^{n_1})\oplus A/(\pi^{n_2})\oplus\cdots,
\qquad n_1,n_2,\ldots\geq 1,
\]
then $M_1$ is not cyclic,
\[
M_1\approx(\pi^{n_1-1})/(\pi^{n_1})\oplus
 (\pi^{n_2-1})/(\pi^{n_2})\oplus\cdots.
\]
Therefore $M_n$ is a cyclic $A$-module of order $q^n$, and every such
module is isomorphic to $A/(\pi^n)$.
\qed
% END CFT_I36_CONTEXT_46
\endgroup

\clearpage
\noindent\emph{Source \texttt{CFT.pdf}, physical page 47.}\par\medskip
% BEGIN CFT_I36_CONTEXT_47
\noindent\textsc{Proposition 3.4}
\emph{The $A$-module $\Lambda_n$ is isomorphic to $A/(\pi^n)$. Hence
$\operatorname{End}_A(\Lambda_n)\simeq A/(\pi^n)$ and
$\operatorname{Aut}_A(\Lambda_n)\simeq(A/(\pi^n))^\times$.}
\par\medskip
\noindent\textsc{Proof.} An $A$-isomorphism $h\colon F_f\longrightarrow F_g$
of formal group laws induces an isomorphism of $A$-modules
$\Lambda_f\longrightarrow\Lambda_g$, and so it does not matter which
$f\in\mathcal F_\pi$ we choose. We take $f\in\mathcal F_\pi$ to be a
polynomial of the form $\pi T+\cdots+T^q$. This is an Eisenstein
polynomial (ANT, p.~129), and so has $q$ distinct roots, each with absolute
value $<1$. Let $\alpha\in K^{\mathrm{al}}$ have absolute value $<1$.
From the Newton polygon of
\[
f(T)-\alpha=T^q+\cdots+\pi T-\alpha
\]
we see that its roots have absolute value $<1$, and so lie in $\Lambda_f$.
We have shown that $\Lambda_f$ satisfies the hypotheses of the lemma, and
so $\Lambda_n\approx A/(\pi^n)$. It follows that the action of $A$ on
$\Lambda_n$ induces an isomorphism
$A/(\pi^n)\longrightarrow\operatorname{End}_A(\Lambda_n)$.
\qed

\medskip
\noindent\textsc{Lemma 3.5}
\emph{Let $L$ be a finite Galois extension of a local field $K$, with
Galois group $G$. For every $F\in\mathcal O_K[[X_1,\ldots,X_n]]$ and
$\alpha_1,\ldots,\alpha_n\in\mathfrak m_L$,}
\[
F(\tau\alpha_1,\ldots,\tau\alpha_n)=\tau F(\alpha_1,\ldots,\alpha_n),
\qquad\text{all $\tau\in G$.}
\]

\noindent\textsc{Proof.} If $F$ is a polynomial, this follows from the
fact that $\tau$ is a field isomorphism fixing the elements of
$\mathcal O_K$. Because the absolute value on $L$ is the unique extension
of the absolute value on $K$ (ANT, 7.38), it is preserved by $\tau\in G$,
i.e., $|\tau\alpha|=|\alpha|$ all $\alpha\in L$, and so $\tau$ is
continuous. Therefore it preserves limits,
\[
\lim_{m\to\infty}\alpha_m=L\quad\Rightarrow\quad
 \lim_{m\to\infty}\tau\alpha_m=\tau L.
\]
Write $F=F_m+\text{terms of degree }\geq m+1$ with $F_m$ a polynomial
of degree $m$. Then
\[
\begin{aligned}
\tau(F(\alpha_1,\ldots,\alpha_n))
 &=\tau\bigl(\lim_{m\to\infty}F_m(\alpha_1,\ldots,\alpha_n)\bigr)\\
 &=\lim_{m\to\infty}\tau(F_m(\alpha_1,\ldots,\alpha_n))\\
 &=\lim_{m\to\infty}F_m(\tau\alpha_1,\ldots,\tau\alpha_n).
\end{aligned}
\]
\qed
% END CFT_I36_CONTEXT_47

\medskip
\noindent\emph{Published typographical erratum for the supporting
Example 3.2. The archive attributes this entry to Raghuram Sundararajan.}
\par\medskip
% BEGIN CFT_I36_CONTEXT_ERRATUM
\noindent (typographical) In Example 3.2 (page 37), the first formula for
$\Lambda_n$ should be $\{\alpha\mid(\alpha+1)^{p^n}=1\}$, rather than
$\{\alpha\mid(\alpha-1)^{p^n}=1\}$.
% END CFT_I36_CONTEXT_ERRATUM
\end{document}
